%%%%%%%%%%%%%%%%%%%%%%%%%%%%%%%%%%%%%%%%%%%%%%%%%%%%%
%             UNITS, EQUATIONS AND TEXT             %
%%%%%%%%%%%%%%%%%%%%%%%%%%%%%%%%%%%%%%%%%%%%%%%%%%%%%

%%Macro for 'where'enviroment has been improved by Andrea :-)

%Units:
\newcommand{\unit}[1]{&& \left[\si{#1}\right]} %\newcommand{\unit}[1]{[\si{#1}]}             <<| Use these if you want equations to be
\newcommand{\unitWh}[1]{[\si{#1}]}             %\newcommand{\eq}[2]{&&\si{#1} &= \si{#2}&&}  <<| centered.. .. will appear scrambled
\newcommand{\numUnit}[1]{\ \si{#1}&}           %                                               | from one equation to the next though..
%Equation:                                     %                                               | and does not work with long equations.. :/
\newcommand{\eq}[2]{\si{#1} &= \si{#2}}
\newcommand{\arw}{&& &\Updownarrow&&}
\newcommand{\eqOne}[2]{\si{#1} &= \si{#2} &\nonumber\\}
\newcommand{\eqTwo}[1]{&\ \ \ \ \si{#1}&}
\newcommand{\eqThree}[1]{&\ \ \ \ \si{#1}&}
%Text:
\newcommand{\tx}[1]{\text{#1}}
%Vectors
\renewcommand{\vec}[1]{\boldsymbol{\mathbf{#1}}}
%Vertical line in equations ie. |_x=y (whereTwo stacks two equalities at the line)



\newcommand{\oleStreg}[1]{ \left.\right\vert\rule{0cm}{.4cm}_{\substack{\rule{0cm}{.15cm}\\ #1 }} }

\newcommand{\oleStregTwo}[2]{ \left.\rule{0cm}{.7cm}\right\vert\rule{0cm}{.7cm}_{\substack{#1 \rule{0cm}{.2cm}\\\vspace{-.1cm}\\ #2}} }
%
%\newenvironment{where}{\indent Where:\\}{}
%\newcommand{\va}[3]{\hspace{16mm} \begin{tabular}{p{40pt} p{360pt} l} {#1} & {#2} & \ifthenelse{\isempty{#3}}{}{[{#3}]} \\
%	\end{tabular}\\}





%%%%%%%%%%%%%%%%%%%%%%%%%%%%%%%%%%%%%%%%%%%%%%%%%%%%%
%              SETTING UP NOTE BOXES                %
%%%%%%%%%%%%%%%%%%%%%%%%%%%%%%%%%%%%%%%%%%%%%%%%%%%%%
%\specialcomment{note}{\[}{\]}
%%{\begin{tcolorbox}
%%                      [grow to right by  = 1.5 cm
%%                      ]}{\vspace{-.3 cm}\end{tcolorbox}}
%
%\newcommand{\toggleNotes}[1]
%{
%  \ifthenelse{\equal{#1}{off}}{\excludecomment{note}}{}
%}

%%%%%%%%%%%%%%%%%%%%%%%%%%%%%%%%%%%%%%%%%%%%%%%%%%%%%
%                 TIKZ SETTINGS                     %
%%%%%%%%%%%%%%%%%%%%%%%%%%%%%%%%%%%%%%%%%%%%%%%%%%%%%
%\usetikzlibrary{arrows.meta}
\tikzset{
  block/.style    = {draw, thick, rectangle,
                     minimum height = 2.1em,
                     minimum width = 1.7em},
  sum/.style      = {draw, circle, inner sep=1.5pt},
}

%%%%%%%%%%%%%%%%%%%%%%%%%%%%%%%%%%%%%%%%%%%%%%%%%%%%%
%                  REFERENCES                       %
%%%%%%%%%%%%%%%%%%%%%%%%%%%%%%%%%%%%%%%%%%%%%%%%%%%%%

%Chapter
\newcommand{\Chapref}[1]{\emph{Chapter \ref{#1}}}
\newcommand{\chapref}[1]{\emph{chapter \ref{#1}}}
%Section
\newcommand{\Secref}[1]{\emph{Section \ref{#1}}}
\newcommand{\secref}[1]{\emph{section \ref{#1}}}
%subSection
\newcommand{\Subsecref}[1]{\emph{Subsection \ref{#1}}}
\newcommand{\subsecref}[1]{\emph{subsection \ref{#1}}}
%Appendix
\newcommand{\Appref}[1]{\emph{Appendix \ref{#1}}}
\newcommand{\appref}[1]{\emph{appendix \ref{#1}}}
%Listings
\newcommand{\Coderef}[1]{\emph{Listings: \ref{#1}}}
\newcommand{\coderef}[1]{\emph{listings: \ref{#1}}}
%Figure:
\newcommand{\Figref}[1]{\emph{Figure \ref{#1}}}
\newcommand{\figref}[1]{\emph{figure \ref{#1}}}
%Table:
\newcommand{\Tableref}[1]{\emph{Table \ref{#1}}}
\newcommand{\tableref}[1]{\emph{table \ref{#1}}}

%Expressions:
\newcommand{\Expr}[1]{\emph{Expression (\ref{#1})}}
\newcommand{\expr}[1]{\emph{expression (\ref{#1})}}

%Equations:
%1 equation:
\newcommand{\Eqref}[1]{\emph{Equation (\ref{#1})}}
\renewcommand{\eqref}[1]{\emph{equation (\ref{#1})}}
%2 equations:
\newcommand{\EqrefTwo}[2]{\emph{Equation (\ref{#1})} and \emph{(\ref{#2})}}
\newcommand{\eqrefTwo}[2]{\emph{equation (\ref{#1})} and \emph{(\ref{#2})}}
%3 equations:
\newcommand{\EqrefThree}[3]{\emph{Equation (\ref{#1})}, \emph{(\ref{#2})} and \emph{(\ref{#3})}}
\newcommand{\eqrefThree}[3]{\emph{equation (\ref{#1})}, \emph{(\ref{#2})} and \emph{(\ref{#3})}}
%4 equations:
\newcommand{\EqrefFour}[4]{\emph{Equation (\ref{#1})}, \emph{(\ref{#2})}, \emph{(\ref{#3})} and \emph{(\ref{#4})}}
\newcommand{\eqrefFour}[4]{\emph{equation (\ref{#1})}, \emph{(\ref{#2})}, \emph{(\ref{#3})} and \emph{(\ref{#4})}}
%5 equations:
\newcommand{\EqrefFive}[5]{\emph{Equation (\ref{#1})}, \emph{(\ref{#2})}, \emph{(\ref{#3})}, \emph{(\ref{#4})} and \emph{(\ref{#5})}}
\newcommand{\eqrefFive}[5]{\emph{equation (\ref{#1})}, \emph{(\ref{#2})}, \emph{(\ref{#3})}, \emph{(\ref{#4})} and \emph{(\ref{#5})}}
%6 equations:
\newcommand{\EqrefSix}[6]{\emph{Equation (\ref{#1})}, \emph{(\ref{#2})}, \emph{(\ref{#3})}, \emph{(\ref{#4})}, \emph{(\ref{#5})} and \emph{(\ref{#6})}}
\newcommand{\eqrefSix}[6]{\emph{equation (\ref{#1})}, \emph{(\ref{#2})}, \emph{(\ref{#3})}, \emph{(\ref{#4})}, \emph{(\ref{#5})} and \emph{(\ref{#6})}}
%7 equations:
\newcommand{\EqrefSeven}[7]{\emph{Equation (\ref{#1})}, \emph{(\ref{#2})}, \emph{(\ref{#3})}, \emph{(\ref{#4})}, \emph{(\ref{#5})}, \emph{(\ref{#6})} and \emph{(\ref{#7})}}
\newcommand{\eqrefSeven}[7]{\emph{equation (\ref{#1})}, \emph{(\ref{#2})}, \emph{(\ref{#3})}, \emph{(\ref{#4})}, \emph{(\ref{#5})}, \emph{(\ref{#6})} and \emph{(\ref{#7})}}


\newcommand{\newpar}{\par \vspace{0.4cm}}


%\usepackage{xifthen}
%
%\newenvironment{where}{\indent Where:\\}{}
%\newcommand{\va}[3]
%{
%  \hspace{-18mm}
%  \begin{tabular}{p{35pt} p{345pt} l}
%    { $#1$ } & { #2 } & \ifthenelse{\isempty{ #3 }}  {}  {\si{[#3]}}
%  \end{tabular}
%}















\usetikzlibrary{calc}
\usetikzlibrary{backgrounds}
\newcommand{\showgrid}{%
  \AtBeginShipoutNext{\AtBeginShipoutAddToBoxForeground{%
      \begin{tikzpicture}
        [
          overlay,
          remember picture,
          inner sep=0pt,
          outer sep=0pt,
          minor line/.style={help lines, draw=gray!25, on background layer},
          major line/.style={help lines, draw=gray},
        ]
        \foreach \step in {0,...,210} {
          \pgfmathsetmacro\gridlineconfig{ifthenelse(equal(int(mod(\step,10)),0),"major line","minor line")}%
          \draw [\gridlineconfig] ($(current page.north west) + (\step mm,0)$) -- ($(current page.south west) + (\step mm,0)$);
        }
        \foreach \step in {0,...,297} {
          \pgfmathsetmacro\gridlineconfig{ifthenelse(equal(int(mod(\step,10)),0),"major line","minor line")}%
          \draw [\gridlineconfig] ($(current page.north west) - (0,\step mm)$) -- ($(current page.north east) - (0,\step mm)$);
          \node [anchor=north] at ($ (current page.north west) + (\step mm,0) $) {\fontsize{1}{2}\selectfont \step};
          \node [anchor=west] at ($ (current page.north west) - (0,\step mm) $) {\fontsize{1}{2}\selectfont \step};
        }
      \end{tikzpicture}
    }%
  }%
}













